\begin{figure}[h]
    \centering
    \newcommand{\radarplot}[6]{ % Define a command for radar plot
        \begin{tikzpicture}
            \begin{polaraxis}[
                ymax=1.1, ymin=0, 
                xtick=\empty, ytick=\empty,
                width=0.3\linewidth,
                title={#1},
                ylabel={#2},
                xticklabels={E,C,A},
            ]
                \addplot+[mark=none,smooth cycle,fill=#6] coordinates {
                    (0,0) (90,#3) (210,#4) (330,#5)
                };
            \end{polaraxis}
        \end{tikzpicture}
    }

    \begin{tabular}{cccc}
             & Problem 1                                & Problem2                              & Problem3                              \\
        SUT1 & \radarplot{}{}{1}{1}{1}{blue!30} & \radarplot{}{}{0.3}{0.7}{0.9}{red!30} & \radarplot{}{}{0.8}{0.5}{0.6}{green!30} \\
        SUT2 & \radarplot{}{}{0.7}{0.9}{0.4}{orange!30} & \radarplot{}{}{0.5}{0.5}{0.5}{purple!30} & \radarplot{}{}{0.4}{0.9}{0.7}{yellow!30} \\
        SUT3 & \radarplot{}{}{0.8}{0.1}{0.9}{cyan!30} & \radarplot{}{}{0.6}{0.5}{0.7}{magenta!30} & \radarplot{}{}{0.9}{0.7}{1.0}{teal!30} \\
    \end{tabular}

    \caption{Mock-Up of Radar Plots to show what the output of a \gai benchmark run looks like. Radar plots are along the axes of effectiveness (north) efficiency (southwest) and cost (southeast)}
    \label{fig:radar_plots}
\end{figure}
